\expandafter\def\csname HistoricalNoteText20\endcsname{Im vorangegangenen Artikel finden sich bereits einige Beispiele für die beträchtliche Abweichung zwischen verschiedenen Beobachtern an ein und demselben Ort. Einige weitere mögen hier angeführt werden, bei denen die Unterschiede viel größer sind, als sich mit irgendeiner Wahrscheinlichkeit durch regelmäßige jährliche Änderung erklären lässt. In Valparaiso betrug die Inklination 1829 nach \textsc{King} $-40^0\,11'$, 1835 nach \textsc{Fitz Roy} $-38^0\,3'$. Auf der Insel Mauritius betrug die Intensität im Jahr 1818 nach \textsc{Freycinet} $1{,}096$, im Jahr 1836 nach \textsc{Fitz Roy} $1{,}192$. Noch größer ist der Unterschied auf Otaheite, wo \textsc{Erman} 1830 die Intensität $=1{,}172$ fand, \textsc{Fitz Roy} dagegen 1835 $=1{,}017$. Dieser letzte Unterschied an einem für die künftige Verbesserung der Elemente höchst wichtigen Ort ist beträchtlich größer als der größte Unterschied in allen unseren 86 Vergleichen berechneter mit beobachteten Intensitäten.}
\SourceLeaf{166}
\Article{30}
Bei der Beurteilung der Unterschiede zwischen Rechnung und Beobachtung, die sich aus dem vorstehenden tabellarischen Vergleich ergeben, muss man zweierlei berücksichtigen. Einerseits sind fast alle Beobachtungen mit Fehlern bei der Messung und mit zufälligen Anomalien der magnetischen Kraft selbst behaftet und gelten nicht für ein und dasselbe Jahr\HistoricalNote{20}. Andererseits enthalten unsere Formeln nur die Glieder bis zur vierten Ordnung, während die folgenden noch sehr merklich sein können. Die Übereinstimmung zwischen Rechnung und Beobachtung scheint daher allen angemessenen Erwartungen an einen ersten Versuch zu genügen. Unser Ausdruck für $\frac VR$ darf also wohl als der Wahrheit nahe betrachtet werden, zumindest in seinen bedeutenderen Gliedern. Es schien daher die Mühe wert, den Verlauf der Zahlenwerte von $\frac VR$ grafisch zu veranschaulichen. Dies ist mit einer von Herrn Dr. \textsc{Goldschmidt} gezeichneten Karte in drei Teilen geschehen. Der erste stellt in \textsc{Mercator}-Projektion den ganzen Erdgürtel zwischen $70^0$ nördlicher und $70^0$ südlicher Breite dar. Die beiden anderen zeigen in stereografischer Projektion die Polargebiete bis zu $65^0$ Breite. Die Korrekturen und Ergänzungen, die eine erneute Rechnung auf der Grundlage besserer Daten künftig am Ausdruck für $\frac VR$ nötig machen wird, werden ohne Zweifel noch beträchtliche Verschiebungen in diesem Liniensystem hervorrufen, besonders in den hohen südlichen Breiten. Eine wesentliche Änderung der Gesamtgestalt ist aber ohne so große Änderungen im Ausdruck für $\frac VR$ nicht denkbar, dass die Übereinstimmung mit den vorhandenen Beobachtungen verloren gehen müsste. Wir gelangen so zu dem wichtigen Ergebnis, dass das System der
\SourceLeaf{167}
Linien gleicher Werte von $V$ auf der Erdoberfläche tatsächlich dem einfachsten oben in Art. 11 beschriebenen Typ entspricht und dass es auf der Erde somit \textit{nur zwei magnetische Pole} gibt, wenn man vom in Artikel 13 erwähnten Fall einer örtlichen Ausnahme absieht. Ob dieser Fall vorkommt, muss vorerst offenbleiben. Die genaue Berechnung nach unseren Elementen ergibt die Lage dieser beiden Pole

1) bei $73^0\,35'$ nördlicher Breite und $264^0\,21'$ Länge östlich von Greenwich mit einer Gesamtintensität $=1{,}701$ (in der üblichen Einheit).

2) bei $72^0\,35'$ südlicher Breite und $152^0\,30'$ Länge mit einer Gesamtintensität $=2{,}253$.

Im ersten Punkt hat $\frac VR$ seinen größten Wert $=+895{,}86$, im zweiten seinen kleinsten $=-1030{,}24$.

Nach \textsc{Ross}’ Beobachtung liegt der nördliche magnetische Pol um $3^0\,30'$ weiter südlich als nach unserer Rechnung. Wie aus unserer Vergleichstafel ersichtlich ist, ergibt die Rechnung für diesen Ort eine um $1^0\,12'$ fehlerhafte Richtung der magnetischen Kraft. Beim südlichen magnetischen Pol ist eine beträchtlich größere Verschiebung zu erwarten. In Hobarttown, dem nächstgelegenen Beobachtungsort, gibt die Rechnung die Inklination ohne Rücksicht auf das Vorzeichen um $3^0\,38'$ zu klein an, soweit man sich auf die Beobachtung verlassen kann. Der wirkliche südliche magnetische Pol liegt daher wahrscheinlich beträchtlich weiter nördlich, als unsere Rechnung angibt, und wäre etwa in der Gegend von $66^0$ Breite und $146^0$ Länge zu suchen.

\Article{31}
Den beiden Punkten auf der Erdoberfläche, an denen die horizontale Kraft verschwindet und die man die magnetischen Pole nennt, mag man wegen ihres Zusammenhangs mit der Gestalt der Erscheinungen der horizontalen Kraft auf der gesamten Erdoberfläche eine gewisse Bedeutung beimessen. Man muss sich aber davor hüten, diese Bedeutung weiter auszudehnen. Insbesondere ist die Sehne, die diese beiden Punkte verbindet, ohne jede Bedeutung, und es wäre verfehlt, \textit{diese} Gerade als \textit{magnetische Achse} der Erde zu bezeichnen. Die einzige Art, dem Begriff der magnetischen Achse eines Körpers eine
\SourceLeaf{168}
allgemein gültige Bedeutung zu geben, ist die in Artikel 5 der \textit{Intensitas vis magneticae} festgelegte. Danach versteht man darunter eine Gerade, bezüglich deren das Moment des im Körper enthaltenen freien Magnetismus ein Maximum ist. Um die Lage der magnetischen Achse der Erde in diesem Sinn und zugleich das Moment des Erdmagnetismus bezüglich dieser Achse zu bestimmen, braucht man, wie oben in Artikel 17 bereits bemerkt, nur die Glieder erster Ordnung von $V$ zu kennen. Nach unseren Elementen in Artikel 26 ist
\[
P'=+925{,}782\cos u+89{,}024\sin u\cos\lambda-178{,}744\sin u\sin\lambda
\]
Somit sind $-925{,}782R^3$, $-89{,}024R^3$, $+178{,}744R^3$ die Momente des Erdmagnetismus bezüglich der Erdachse und der beiden Erdradien für die Längen $0$ und $90^0$. Bei der Erdachse ist die Richtung zum Nordpol gemeint. Das negative Vorzeichen des entsprechenden Moments zeigt an, dass die magnetische Achse mit ihr einen stumpfen Winkel bildet, das heißt, dass ihr magnetischer Nordpol nach Süden gekehrt ist. Daraus ergibt sich die Richtung der magnetischen Achse parallel zu dem Erddurchmesser, der von $77^0\,50'$ nördlicher Breite und $296^0\,29'$ Länge nach $77^0\,50'$ südlicher Breite und $116^0\,29'$ Länge verläuft. Das magnetische Moment bezüglich dieser Achse ist $=947{,}08R^3$. Dabei ist zu bedenken, dass unseren Elementen eine Intensitätseinheit zugrunde liegt, die ein Tausendstel der üblichen Einheit ist. Für die Umrechnung auf die in der \textit{Intensitas vis magneticae} festgelegte absolute Einheit ist zu beachten, dass dort die horizontale Intensität in Göttingen am 19. Juli 1834 als $=1{,}7748$ bestimmt wurde. Daraus folgt mit der Inklination $68^0\,1'$ die Gesamtintensität $=4{,}7414$, während sie in der oben verwendeten Einheit mit $=1357$ angesetzt wird. Der Umrechnungsfaktor ist also $=0{,}0034941$, und damit ist das magnetische Moment der Erde in der absoluten Einheit
\[
=3{,}3092R^3
\]
Da bei dieser absoluten Einheit für die erdmagnetische Kraft der Millimeter als Längeneinheit gilt, muss auch $R$ in Millimetern angesetzt werden. Da die Elliptizität der Erde hier ohnehin unberücksichtigt bleibt, genügt es, $R$ als Radius eines Kreises mit einem Umfang von $40000$ Millionen Millimetern zu betrachten. Das obige magnetische Moment wird demnach durch
\expandafter\def\csname HistoricalNoteText21\endcsname{Insofern wir nämlich nicht berechtigt sind, bei allen magnetisierten Teilen der Erde durchweg parallele magnetische Achsen vorauszusetzen. Je stärker die Achsen vom Parallelismus abweichen, desto stärker muss die durchschnittliche Magnetisierung der Teile sein, damit dasselbe magnetische Gesamtmoment entsteht.}
\SourceLeaf{169}
eine Zahl ausgedrückt, deren Logarithmus $=29{,}93136$ ist, oder durch $853800$ Quadrillionen. In derselben absoluten Einheit wurde das magnetische Moment eines einpfündigen Magnetstabs in den Versuchen von 1832 als $=100877000$ bestimmt (\textit{Intensitas} Art. 21). Das magnetische Moment der Erde ist also $8464$ Trillionen Mal so groß. Es wären daher $8464$ Trillionen solcher Magnetstäbe mit parallelen magnetischen Achsen erforderlich, um die magnetische Wirkung der Erde im Außenraum zu ersetzen. Bei gleichmäßiger Verteilung über das gesamte Volumen der Erde entspricht das beinahe acht Stäben (genauer $7{,}831$) pro Kubikmeter. So ausgedrückt behält das Ergebnis seine Bedeutung auch dann, wenn man die Erde nicht als wirklichen Magneten auffasst, sondern den Erdmagnetismus allein auf beständige galvanische Ströme in der Erde zurückführt. Betrachten wir die Erde aber als wirklichen Magneten, so müssen wir jedem ihrer Teile von einem Achtel Kubikmeter \textit{im Durchschnitt} wenigstens\HistoricalNote{21} eine ebenso starke Magnetisierung zuschreiben, wie sie jener Magnetstab besitzt. Dieses Ergebnis dürfte die Physiker überraschen.

